% Quelle: 5. Übung CT Lösungen.pdf
% Art: Übung mit Lösungen
% Themen: Lineare Codes, Wiederholungscode, Paritätscode, Erzeugermatrix, Kontrollmatrix, Syndrom, Hamming-Code H_2(m), Hamming-Code H_q(m), Dekodieren, Hamming-Schranke, Singleton-Schranke, perfekte Codes, optimale Codes
\section{Übung 5 (mit Lösungen)}
% Seite 1: Aufgabenblatt, Seiten 2--5: "Ü5 Lösungen".
% In den Matrizen der Quelle sind Nulleinträge meist leer gelassen; leere Zellen wurden so übernommen (leer = 0).

\subsection{Wiederholungscode $C \subseteq \mathbf{F}_3^5$}

\begin{aufgabe}[5\_1]
\begin{enumerate}[label=\alph*)]
\item Geben Sie eine Erzeugermatrix $G$ und eine Kontrollmatrix $H$ für $C$ an.
\item Codieren Sie das Wort $a = 2 \in \mathbf{F}_3^1$.
\item Prüfen Sie, ob $\mathbf{y}_1 = 22120$, $\mathbf{y}_2 = 11022$, $\mathbf{y}_3 = 11111 \in \mathbf{F}_3^5$ Codewörter sind.
\end{enumerate}
\end{aufgabe}

\begin{loesung}
\begin{enumerate}[label=\alph*)]
\item
\[
G = (\,1 \mid 1\;1\;1\;1\,) = (\,E \mid A\,), \qquad
H = (-A^T \mid E) =
\begin{pmatrix}
2 & 1 & & & \\
2 & & 1 & & \\
2 & & & 1 & \\
2 & & & & 1
\end{pmatrix}
\]
\item $a\,G = 2\;2\;2\;2\;2$
\item
\begin{align*}
S(\mathbf{y}_1) &= H\,\mathbf{y}_1^T = (0\;2\;0\;1)^T \ne \mathbf{o}\\
S(\mathbf{y}_2) &= H\,\mathbf{y}_2^T = (0\;2\;1\;1)^T \ne \mathbf{o}\\
S(\mathbf{y}_3) &= H\,\mathbf{y}_3^T = (0\;0\;0\;0)^T = \mathbf{o} \;\Rightarrow\; \mathbf{y}_3 \in C
\end{align*}
\end{enumerate}
\end{loesung}

\subsection{Paritätscode $C \subseteq \mathbf{F}_3^5$}

\begin{aufgabe}[5\_2]
\begin{enumerate}[label=\alph*)]
\item Geben Sie eine Erzeugermatrix $G$ und eine Kontrollmatrix $H$ für $C$ an.
\item Codieren Sie das Wort $\mathbf{a} = 2120 \in \mathbf{F}_3^4$.
\item Prüfen Sie, ob $\mathbf{y}_1 = 22120$, $\mathbf{y}_2 = 11022$, $\mathbf{y}_3 = 11111 \in \mathbf{F}_3^5$ Codewörter sind.
\end{enumerate}
\end{aufgabe}

\begin{loesung}
\begin{enumerate}[label=\alph*)]
\item
\[
G = (\,E \mid A\,) =
\begin{pmatrix}
1 & & & & 1\\
& 1 & & & 1\\
& & 1 & & 1\\
& & & 1 & 1
\end{pmatrix}, \qquad
H = (-A^T \mid E) = (\,2\;2\;2\;2 \mid 1\,)
\]
% KORRIGIERT: "a G = 2 1 2 0 G" (a nicht fett, Klammern fehlten)
\item $\mathbf{a}\,G = (\,2\;1\;2\;0\,)\,G = 2\;1\;2\;0\;2$ \quad (Paritätsstelle $2+1+2+0 = 5 \equiv 2 \pmod{3}$)
\item
% KORRIGIERT: Syndrom ist hier ein Skalar in F_3; "\ne \mathbf{o}" -> "\ne 0" (einheitlich mit der letzten Zeile)
\begin{align*}
S(\mathbf{y}_1) &= H\,\mathbf{y}_1^T = (\,2\;2\;2\;2\;\;1\,)\;\mathbf{y}_1^T = 2 \ne 0\\
S(\mathbf{y}_2) &= H\,\mathbf{y}_2^T = (\,2\;2\;2\;2\;\;1\,)\;\mathbf{y}_2^T = 1 \ne 0\\
S(\mathbf{y}_3) &= H\,\mathbf{y}_3^T = (\,2\;2\;2\;2\;\;1\,)\;\mathbf{y}_3^T = 0 \;\Rightarrow\; \mathbf{y}_3 \in C
\end{align*}
\end{enumerate}
\end{loesung}

\subsection{Hamming Code $H_2(m)$}

\begin{aufgabe}[5\_3]
\begin{enumerate}[label=\alph*)]
\item Bestimmen Sie eine Kontrollmatrix $H$ und eine Erzeugermatrix $G$ für den Hamming Code $H(2,2)$.
\item Bestimmen Sie eine Kontrollmatrix $H$ für den Hamming Code $H_2(4)$.
\end{enumerate}
% Auf dem Aufgabenblatt steht in a) "H(2, 2)", in der Lösung "H_2(2)".
\end{aufgabe}

\begin{loesung}
Bestimmen Sie eine Kontrollmatrix $H$ und eine Erzeugermatrix $G$ für den
\begin{enumerate}[label=\alph*)]
\item Hamming Code $H_2(2)$

Kontrollmatrix
\[
H = \begin{pmatrix} 1 & 0 & 1\\ 0 & 1 & 1 \end{pmatrix}, \qquad
G = (-A^T \mid E) = (\,1\;\;1\;\;1\,)
\]
% KORRIGIERT: "F_2^3" (nicht fett, uneinheitlich)
Wiederholungscode $C \subseteq \mathbf{F}_2^3$
\item Hamming Code $H_2(4)$
\[
H =
\begin{pmatrix}
1 & & & & & & & & 1 & 1 & 1 & 1 & 1 & 1 & 1\\
& 1 & & & & 1 & 1 & 1 & & & & 1 & 1 & 1 & 1\\
& & 1 & & 1 & & 1 & 1 & & 1 & 1 & & & 1 & 1\\
& & & 1 & 1 & 1 & & 1 & 1 & 0 & 1 & & 1 & & 1
\end{pmatrix}
\]
% In Zeile 4, Spalte 10 steht in der Quelle explizit eine 0 (sonst sind Nullen leer gelassen).
\end{enumerate}
\end{loesung}

\subsection{Dekodieren in $H_2(3)$}

\begin{aufgabe}[5\_4]
Dekodieren Sie die empfangenen Wörter
\[
\mathbf{y}_1 = 1\;1\;0\;1\;1\;0\;0 \qquad
\mathbf{y}_2 = 1\;1\;1\;1\;1\;1\;1 \qquad
\mathbf{y}_3 = 1\;1\;1\;1\;0\;0\;0
\]
\end{aufgabe}

\begin{loesung}
\[
H =
\begin{pmatrix}
1 & & & 1 & 1 & 0 & 1\\
& 1 & & 1 & 0 & 1 & 1\\
& & 1 & 0 & 1 & 1 & 1
\end{pmatrix}
\]
\begin{align*}
S(\mathbf{y}_1) &= H\,\mathbf{y}_1^T = (\,1\;0\;1\,)^T = \mathbf{h}_5 && \Rightarrow\; \mathbf{c}_1 = 1\;1\;0\;1\;0\;0\;0\\
S(\mathbf{y}_2) &= H\,\mathbf{y}_2^T = (\,0\;0\;0\,)^T = \mathbf{o} && \Rightarrow\; \mathbf{c}_2 = \mathbf{y}_2\\
S(\mathbf{y}_3) &= H\,\mathbf{y}_3^T = (\,0\;0\;1\,)^T = \mathbf{h}_3 && \Rightarrow\; \mathbf{c}_3 = 1\;1\;0\;1\;0\;0\;0
\end{align*}
\end{loesung}

\subsection{Hamming Code $H_q(m)$}

\begin{aufgabe}[5\_5]
Die Spalten einer Kontrollmatrix $H$ des $q$-nären Codes $H_q(m)$ bestehen aus Repräsentanten der 1-dimensionalen Unterräume des $\mathbf{F}_q^m$.
\begin{enumerate}[label=\alph*)]
\item Wie viele Vektoren besitzt ein 1-dimensionaler Unterraum des $\mathbf{F}_q^n$?
\item Wie lauten die Parameter $[n, k, d]$ des Codes $H_q(m)$?
\item Bestimmen Sie eine Kontrollmatrix $H$ für den Hamming Code $H_3(3)$.
\item Dekodieren Sie die Wörter
\begin{align*}
\mathbf{r}_1 &= 0\,0\,0\,0\,0\,0\,0\,0\,0\,0\,1\,1\,2\\
\mathbf{r}_2 &= 0\,0\,0\,0\,0\,0\,0\,0\,0\,0\,1\,2\,2
\end{align*}
\end{enumerate}
% Aufgabenblatt: in a) steht "F_q^n", in der Lösung "F_q^m" und "Elemente" statt "Vektoren".
\end{aufgabe}

\begin{loesung}
Die Spalten einer Kontrollmatrix $H$ des $q$-nären Hamming Codes $H_q(m)$ bestehen aus Repräsentanten der 1-dimensionalen Unterräume des $\mathbf{F}_q^m$.
\begin{enumerate}[label=\alph*)]
\item Wie viele Elemente besitzt ein 1-dimensionaler Unterraum des $\mathbf{F}_q^m$?

% KORRIGIERT: Voraussetzung an h fehlte
1-dim.\ Unterraum des $\mathbf{F}_q^m$ sind die Geraden (mit $\mathbf{h} \in \mathbf{F}_q^m$, $\mathbf{h} \ne \mathbf{o}$)
\[
\{\, y\,\mathbf{h} \mid y \in \mathbf{F}_q \,\}, \qquad |\{\, y\,\mathbf{h} \mid y \in \mathbf{F}_q \,\}| = |\mathbf{F}_q| = q
\]
\item Wie lauten die Parameter $[n, k, d]$ des Codes $H_q(m)$?
\begin{align*}
& \operatorname{Rang} G + \operatorname{Rang} H = n = (q^m - 1)/(q-1)\\
& \operatorname{Rang} H = m, \quad \operatorname{Rang} G = k = n - m\\
\Rightarrow\quad & |C| = q^k
\end{align*}
\item Bestimmen Sie eine Kontrollmatrix $H$ für den Hamming Code $H_3(3)$.
\begin{align*}
n &= (3^3 - 1)/(3-1) = 26/2 = 13\\
|H_3(3)| &= 3^{n-m} = 3^{10}
\end{align*}
\[
H =
\begin{pmatrix}
1 & 0 & 0 & 0 & 0 & 1 & 1 & 1 & 1 & 1 & 1 & 1 & 1\\
0 & 1 & 0 & 1 & 1 & 0 & 0 & 1 & 1 & 1 & 2 & 2 & 2\\
0 & 0 & 1 & 1 & 2 & 1 & 2 & 0 & 1 & 2 & 0 & 1 & 2
\end{pmatrix}
\]
\item Dekodieren Sie die Wörter
\begin{align*}
\mathbf{r}_1 &= 0\,0\,0\,0\,0\,0\,0\,0\,0\,0\,1\,1\,2\\
\mathbf{r}_2 &= 0\,0\,0\,0\,0\,0\,0\,0\,0\,0\,1\,2\,2
\end{align*}
\begin{align*}
% KORRIGIERT: "S(r_1) = ..." ohne H r_1^T; Rechnung in F_3 (mod 3)
S(\mathbf{r}_1) &= H\,\mathbf{r}_1^T = \begin{pmatrix} 1\\ 2\\ 2 \end{pmatrix} = \mathbf{h}_{13}
\;\Rightarrow\; x = 13,\ y = 1 \;\Rightarrow\; \mathbf{c}_1 = 0\,0\,0\,0\,0\,0\,0\,0\,0\,0\,1\,1\,1\\
% KORRIGIERT: in der Quelle stand x = 13; (1 2 0)^T ist Spalte h_11, also x = 11
S(\mathbf{r}_2) &= H\,\mathbf{r}_2^T = \begin{pmatrix} 2\\ 1\\ 0 \end{pmatrix} = 2\begin{pmatrix} 1\\ 2\\ 0 \end{pmatrix} = 2\,\mathbf{h}_{11}
\;\Rightarrow\; x = 11,\ y = 2 \;\Rightarrow\; \mathbf{c}_2 = 0\,0\,0\,0\,0\,0\,0\,0\,0\,0\,2\,2\,2
\end{align*}
\end{enumerate}
\end{loesung}

\subsection{obere Schranken}

\begin{aufgabe}[5\_6]
Prüfen Sie die Hamming Schranke für lineare Codes mit den Parametern $[23, 12, 7]_2$, $[11, 6, 5]_3$.
\end{aufgabe}

\begin{loesung}
% KORRIGIERT: "|F_q^n|" (nicht fett); e = Anzahl korrigierbarer Fehler ergänzt
Hamming Schranke: \quad $|C| \displaystyle\sum_{i=0}^{e} \binom{n}{i} (q-1)^i \le |\mathbf{F}_q^n| = q^n$, \quad $e = \lfloor (d-1)/2 \rfloor$

$[23, 12, 7]_2$
\begin{align*}
% KORRIGIERT: "C = 2^k" statt "|C| = 2^k"
& e = 3,\quad |C| = 2^k = 2^{12} = 4096,\quad
\sum_{i=0}^{e} \binom{n}{i} (q-1)^i = 1 + 23 + \frac{23\cdot 22}{1\cdot 2} + \frac{23\cdot 22\cdot 21}{1\cdot 2\cdot 3}\\
&= 24 + 253 + 1771 = 2048 = 2^{11}
\end{align*}
Tatsächlich ist
\[
|C| \sum_{i=0}^{e} \binom{n}{i} (q-1)^i = 2^{12}\cdot 2^{11} = 2^{23} = |\mathbf{F}_2^{23}|
\]

$[11, 6, 5]_3$
\begin{align*}
% KORRIGIERT: "C = 3^k" statt "|C| = 3^k"
& e = 2,\quad |C| = 3^k = 3^6 = 729,\quad
\sum_{i=0}^{e} \binom{n}{i} (q-1)^i = 1 + 11\cdot 2 + \frac{11\cdot 10}{1\cdot 2}\cdot 2^2\\
&= 23 + 220 = 243 = 3^5
\end{align*}
Tatsächlich ist
\[
|C| \sum_{i=0}^{e} \binom{n}{i} (q-1)^i = 3^{6}\cdot 3^{5} = 3^{11} = |\mathbf{F}_3^{11}|
\]
\end{loesung}

\begin{aufgabe}[5\_7]
Wie hoch muss die Wortlänge $n$ eines linearen $[n, 5, 5]_2$ Codes mindestens sein?
\end{aufgabe}

\begin{loesung}
$[n, 5, 5]_2$

Singleton Schranke: \quad $k + d \le n + 1 \;\Rightarrow\; 5 + 5 \le n + 1 \;\Rightarrow\; 9 \le n$

Hamming Schranke: \quad $|C| \displaystyle\sum_{i=0}^{e} \binom{n}{i} (q-1)^i \le |\mathbf{F}_q^n|$, \quad $2e + 1 = d$
\begin{align*}
% KORRIGIERT: e = 2 (aus d = 5) ergänzt
\Rightarrow\quad & 2^k\,(\,1 + n + n\,(n-1)/2\,) \le 2^n, \quad q = 2,\ k = 5,\ e = 2\\
\Rightarrow\quad & 1 + n + n\,(n-1)/2 \le 2^{n-5}
\end{align*}
% KORRIGIERT: für n = 9, 10, 11 stand "\le"; es gilt 46 > 16, 56 > 32, 67 > 64 (Schranke verletzt)
\begin{tabular}{lll}
$n = 9$ & $\Rightarrow\; 1 + 9 + 36 = 46$ & $> 2^4 = 16$\\
$n = 10$ & $\Rightarrow\; 1 + 10 + 45 = 56$ & $> 2^5 = 32$\\
$n = 11$ & $\Rightarrow\; 1 + 11 + 55 = 67$ & $> 2^6 = 64$\\
$n = 12$ & $\Rightarrow\; 1 + 12 + 66 = 79$ & $\le 2^7 = 128$
\end{tabular}
% Die Quelle endet hier ohne ausformuliertes Ergebnis (nicht ergänzt).
\end{loesung}

\begin{aufgabe}[5\_8]
Prüfen Sie für die linearen Codes mit den Parametern $[10, 8, 3]_{11}$, $[10, 6, 5]_{11}$ sowohl die Hamming als auch die Singleton Schranke. Sind diese Codes perfekt oder optimal?
\end{aufgabe}

\begin{loesung}
$[10, 8, 3]_{11}$

% KORRIGIERT: "|F_q^n|" (nicht fett)
Hamming Schranke: \quad $|C| \displaystyle\sum_{i=0}^{e} \binom{n}{i} (q-1)^i \le |\mathbf{F}_q^n|$
\begin{align*}
% KORRIGIERT: "C = 11^8" statt "|C| = 11^8"
& e = 1,\quad k = 8,\quad |C| = 11^8,\quad \sum_{i=0}^{e} \binom{n}{i} (q-1)^i = 1 + 10\cdot 10 = 101\\
& |\mathbf{F}_{11}^{10}| = 11^8\cdot 11^2 = 11^8\cdot 121
\end{align*}
% KORRIGIERT: "zu einem Codewort" -> "zu jedem Codewort" (Quantor)
d.h.\ es gibt $11^8\cdot 121 - 11^8\cdot 101 = 11^8\cdot 20$ Wörter in $\mathbf{F}_{11}^{10}$, deren Abstand zu jedem Codewort mindestens 2 beträgt.

Singleton Schranke \quad $k + d \le n + 1$\\
% KORRIGIERT: Vergleich fehlte
$k = 8,\ d = 3,\quad n = 10$: \quad $k + d = 11 = n + 1$\\
d.h.\ der Code mit den Parametern $[10, 8, 3]_{11}$ ist optimal.

$[10, 6, 5]_{11}$
\begin{align*}
% KORRIGIERT: "C = 11^6" statt "|C| = 11^6"
& e = 2,\quad k = 6,\quad |C| = 11^6,\quad \sum_{i=0}^{e} \binom{n}{i} (q-1)^i = 1 + 10\cdot 10 + 45\cdot 100 = 4601\\
& |\mathbf{F}_{11}^{10}| = 11^6\cdot 11^4 = 11^6\cdot 14641
\end{align*}
% KORRIGIERT: "zu einem Codewort" -> "zu jedem Codewort" (Quantor)
d.h.\ es gibt $11^6\cdot 14641 - 11^6\cdot 4601 = 11^6\cdot 10040$ Wörter in $\mathbf{F}_{11}^{10}$, deren Abstand zu jedem Codewort mindestens 3 beträgt.

Singleton Schranke \quad $k + d \le n + 1$\\
% KORRIGIERT: Vergleich fehlte; im letzten Satz stand [10, 8, 3]_11
$k = 6,\ d = 5,\quad n = 10$: \quad $k + d = 11 = n + 1$\\
d.h.\ der Code mit den Parametern $[10, 6, 5]_{11}$ ist optimal.
\end{loesung}
